\documentclass[11pt]{article}
\usepackage[margin=1in]{geometry}
\usepackage{graphicx}

\title{CSE 351: Artificial Intelligence -- Assignment 02}
\author{mb09640, ha10609}
\date{Fall 2026}

\begin{document}
\maketitle

\section*{Question 1: Neural Networks as Universal Approximators}

\subsection*{Baseline Model Results}
A Multi-Layer Feedforward Neural Network (MFNN) was trained using backpropagation for 5000 epochs with 1 hidden layer (40 neurons), \texttt{tansig} activation, learning rate of 0.05, and initial bias of 0.5 to approximate $f(x, y) = e^{-\frac{x^2+y^2}{8}}\cos(\sqrt{x^2+y^2}+c)$ with $c = 1 \times 10^{-3}$.

\begin{table}[h!]
\centering
\begin{tabular}{|l|c|}
\hline
\textbf{Metric} & \textbf{Value} \tabularnewline \hline
Initial Cost $J$ (Epoch 0) & 0.976229 \tabularnewline \hline
Final Cost $J$ (Epoch 5000) & 0.000185 \tabularnewline \hline
Training MSE & 0.000370 \tabularnewline \hline
Testing MSE & 0.000394 \tabularnewline \hline
\end{tabular}
\caption{Baseline MFNN Training and Testing Results}
\end{table}

% Uncomment the lines below to include your saved plots:
% \begin{center}
% \includegraphics[width=0.45\textwidth]{cost_vs_epoch.png}
% \includegraphics[width=0.45\textwidth]{pred_vs_true.png}
% \end{center}

\subsection*{Fine-Tuning and Hyperparameter Findings}

\begin{table}[h!]
\centering
\begin{tabular}{|l|l|c|c|}
\hline
\textbf{Hyperparameter} & \textbf{Configuration} & \textbf{Train MSE} & \textbf{Test MSE} \tabularnewline \hline
Learning Rate & lr = 0.01 & 0.003373 & 0.003497 \tabularnewline
              & lr = 0.05 (Baseline) & 0.000370 & 0.000394 \tabularnewline
              & lr = 0.1 & 0.000116 & 0.000124 \tabularnewline
              & lr = 0.3 & 0.000127 & 0.000150 \tabularnewline \hline
Number of Layers & 1 Layer: [40] (Baseline) & 0.000370 & 0.000394 \tabularnewline
                 & 2 Layers: [40, 20] & 0.000051 & 0.000060 \tabularnewline
                 & 3 Layers: [40, 20, 10] & 0.000038 & 0.000049 \tabularnewline \hline
Activation Function & tansig (Baseline) & 0.000370 & 0.000394 \tabularnewline
                    & logsig & 0.004096 & 0.004191 \tabularnewline
                    & relu & 0.000284 & 0.000462 \tabularnewline \hline
Initial Bias & bias = 0.0 & 0.001614 & 0.001686 \tabularnewline
             & bias = 0.5 (Baseline) & 0.000370 & 0.000394 \tabularnewline
             & bias = 1.0 & 0.000063 & 0.000068 \tabularnewline
             & bias = 2.0 & 0.000073 & 0.000090 \tabularnewline \hline
\end{tabular}
\caption{Effect of Hyperparameters on Train and Test MSE}
\end{table}

\noindent \textbf{Findings:}
\begin{itemize}
    \item \textbf{Number of Layers:} A single hidden layer of 40 neurons accurately approximates the function, while adding 2 and 3 hidden layers further reduces Test MSE from 0.000394 down to 0.000049 due to higher model capacity.
    \item \textbf{Activation Function:} \texttt{tansig} achieves the best generalization (Test MSE: 0.000394) because it is smooth and zero-centered. \texttt{logsig} converges much slower (Test MSE: 0.004191) due to smaller gradient derivatives, while \texttt{relu} produces piecewise-linear segments that slightly overfit the training points (lower Train MSE of 0.000284, but worse Test MSE of 0.000462).
    \item \textbf{Biases:} Setting initial bias to 0.0 gives higher error (Test MSE: 0.001686) because all hidden neurons output zero at the origin $(0,0)$ where the target function peaks at 1.0. Non-zero biases (0.5 and 1.0) break symmetry at the origin and lower Test MSE to 0.000068, though pushing bias too high (2.0) begins to saturate the neurons.
    \item \textbf{Learning Rates:} A small learning rate (0.01) converges too slowly within 5000 epochs (Test MSE: 0.003497). Increasing the rate to 0.1 yields the lowest error (Test MSE: 0.000124), whereas 0.3 introduces gradient oscillations that slightly degrade final accuracy (Test MSE: 0.000150).
\end{itemize}

\section*{Question 2: Decision Tree Induction}

For this part we wrote our own decision tree in Python. At every node we compute the weighted Gini for each leftover attribute and split on the one with the smallest value. Smaller weighted Gini means the groups after the split are cleaner. The three datasets are in \texttt{DATASET.mat}.

\subsection*{a) GINI Index}

\begin{table}[h!]
\centering
\begin{tabular}{|l|c|c|c|c|c|c|}
\hline
\textbf{Dataset} & $a_1$ & $a_2$ & $a_3$ & $a_4$ & $a_5$ & $a_6$ \tabularnewline \hline
Data1 & \textbf{0.1863} & 0.4935 & 0.4949 & 0.4620 & 0.4861 & 0.4664 \tabularnewline \hline
Data2 & 0.4928 & 0.4899 & 0.4885 & \textbf{0.4834} & 0.4900 & 0.4940 \tabularnewline \hline
Data3 & 0.4690 & 0.4639 & 0.4725 & \textbf{0.4447} & 0.4685 & 0.4691 \tabularnewline \hline
\end{tabular}
\caption{Weighted GINI on the training sets. Bold means best attribute for the root.}
\end{table}

So the root should be $a_1$ for Data1, and $a_4$ for both Data2 and Data3.

Data1 is the easy one. $a_1$ jumps out right away (0.1863). Everything else is still around 0.46 to 0.49. That fits what the prompt says about Data1 being based on simple equality and OR conditions.

Data2 looks basically flat. All six numbers are near 0.49, and even the winner ($a_4$ at 0.4834) is barely better than the others. Since Data2 is parity based, the label depends on a pattern across attributes, not on one column by itself. Gini only checks one attribute at a time, so it does not have much to work with here.

Data3 is in between. $a_4$ wins again (0.4447), so Gini still finds something useful, but the gap is smaller than on Data1. The prompt says Data3 has harder rules and 5\% noise. In this kind of synthetic data that usually means about 5\% of the class labels are flipped on purpose, so a few rows are marked with the wrong class. The attributes stay the same. That makes the training data a bit messy and can hurt how clean the Gini scores look.

Overall, Gini is fine when one attribute actually carries the answer (Data1). It is pretty weak on Data2. On Data3 it still picks a reasonable root, but the noise and harder rules make the scores less separated.

\subsection*{b) Tree Construction (Data1)}

We only had to build the full tree for Data1. Same Gini rule as above: keep splitting until a branch is pure, or until every remaining attribute is constant. Then we ran the tree on the Data1 test set and plotted it with NetworkX.

\begin{table}[h!]
\centering
\begin{tabular}{|l|c|}
\hline
\textbf{Metric} & \textbf{Value} \tabularnewline \hline
Root attribute & $a_1$ \tabularnewline \hline
Data1 test accuracy & 1.0000 (100\%) \tabularnewline \hline
\end{tabular}
\caption{Data1 tree results}
\end{table}

Test accuracy was 100\%. The tree starts at $a_1$, branches for values 1, 2, and 3, and ends in class 0 or class 1 leaves. For Data1 that is what you would hope for, since the rules are simple and there is no label noise.

\section*{Question 3: Image Classification for Object Recognition}

We used the CIFAR-10 files in \texttt{cifar-10-batches-py}. Each image is $32 \times 32$ with 3 color channels, so after flattening we get 3072 features. We divided pixels by 255 and then used \texttt{StandardScaler} (fit on train only).

We did not train on all 50,000 images for the main table. Linear SVM and kNN were taking too long on the full set, so we randomly picked 10,000 training images with seed 42. Testing still used all 10,000 test images. The assignment says this is okay if we mention it.

Settings we used:
\begin{itemize}
    \item Decision Tree: \texttt{max\_depth=15}
    \item ANN: one hidden layer with 80 neurons, \texttt{max\_iter=30}, early stopping on
    \item SVM: linear kernel
    \item kNN: $k=3$
\end{itemize}

\begin{table}[h!]
\centering
\begin{tabular}{|l|c|c|c|c|}
\hline
\textbf{Classifier} & \textbf{Accuracy} & \textbf{Precision} & \textbf{Recall} & \textbf{F-measure} \tabularnewline \hline
(a) Decision Tree & 0.2445 & 0.2438 & 0.2445 & 0.2434 \tabularnewline \hline
(b) Artificial Neural Network & 0.4093 & 0.4094 & 0.4093 & 0.4090 \tabularnewline \hline
(c) Support Vector Machine (SVM) & 0.3026 & 0.3086 & 0.3026 & 0.3032 \tabularnewline \hline
(d) K-Nearest Neighbour (kNN) & 0.2850 & 0.3820 & 0.2850 & 0.2709 \tabularnewline \hline
\end{tabular}
\caption{CIFAR-10 scores with 10,000 training images and the full test set. Precision, recall, and F-measure are macro averages.}
\end{table}

ANN wins by a decent margin (about 41\%). Then SVM around 30\%, kNN around 29\%, and the tree around 24\%. Chance is 10\%, so none of them are useless, but none of them are great either. Using raw pixels for this dataset is hard.

The ANN can mix the pixel values together, so it picks up more than a straight linear SVM. The tree is mostly looking at individual pixels / coarse colors, which is why it keeps confusing similar classes (plane vs ship, car vs truck, cat vs dog). kNN has a related issue: two pictures can look close in RGB space and still be different objects. Its precision is higher than its accuracy, so it seems stronger on some classes than others.

We also tried other hyperparameters on the same 10k split (tree depth 20, ANN with 100 units and 20 iterations, $k=5$). Accuracies were 0.2472, 0.4173, 0.3026, and 0.2934. A bit better in places, but nothing dramatic. Changing these settings a little does not fix the raw-pixel problem.

Before we switched to 10k, we did get the tree and ANN to finish on the full 50k training set. Those accuracies were 0.2755 and 0.4985. So more data helps, especially for the ANN. We stopped the full run because the linear SVM was going to take too long, which is why the comparison table above uses the 10k subsample.

\end{document}
